% Options for packages loaded elsewhere
\PassOptionsToPackage{unicode}{hyperref}
\PassOptionsToPackage{hyphens}{url}
\documentclass[
  ignorenonframetext,
]{beamer}
\newif\ifbibliography
\usepackage{pgfpages}
\setbeamertemplate{caption}[numbered]
\setbeamertemplate{caption label separator}{: }
\setbeamercolor{caption name}{fg=normal text.fg}
\beamertemplatenavigationsymbolsempty
% remove section numbering
\setbeamertemplate{part page}{
  \centering
  \begin{beamercolorbox}[sep=16pt,center]{part title}
    \usebeamerfont{part title}\insertpart\par
  \end{beamercolorbox}
}
\setbeamertemplate{section page}{
  \centering
  \begin{beamercolorbox}[sep=12pt,center]{section title}
    \usebeamerfont{section title}\insertsection\par
  \end{beamercolorbox}
}
\setbeamertemplate{subsection page}{
  \centering
  \begin{beamercolorbox}[sep=8pt,center]{subsection title}
    \usebeamerfont{subsection title}\insertsubsection\par
  \end{beamercolorbox}
}
% Prevent slide breaks in the middle of a paragraph
\widowpenalties 1 10000
\raggedbottom
\AtBeginPart{
  \frame{\partpage}
}
\AtBeginSection{
  \ifbibliography
  \else
    \frame{\sectionpage}
  \fi
}
\AtBeginSubsection{
  \frame{\subsectionpage}
}
\usepackage{iftex}
\ifPDFTeX
  \usepackage[T1]{fontenc}
  \usepackage[utf8]{inputenc}
  \usepackage{textcomp} % provide euro and other symbols
\else % if luatex or xetex
  \usepackage{unicode-math} % this also loads fontspec
  \defaultfontfeatures{Scale=MatchLowercase}
  \defaultfontfeatures[\rmfamily]{Ligatures=TeX,Scale=1}
\fi
\usepackage{lmodern}
\ifPDFTeX\else
  % xetex/luatex font selection
\fi
% Use upquote if available, for straight quotes in verbatim environments
\IfFileExists{upquote.sty}{\usepackage{upquote}}{}
\IfFileExists{microtype.sty}{% use microtype if available
  \usepackage[]{microtype}
  \UseMicrotypeSet[protrusion]{basicmath} % disable protrusion for tt fonts
}{}
\makeatletter
\@ifundefined{KOMAClassName}{% if non-KOMA class
  \IfFileExists{parskip.sty}{%
    \usepackage{parskip}
  }{% else
    \setlength{\parindent}{0pt}
    \setlength{\parskip}{6pt plus 2pt minus 1pt}}
}{% if KOMA class
  \KOMAoptions{parskip=half}}
\makeatother
\usepackage{graphicx}
\makeatletter
\newsavebox\pandoc@box
\newcommand*\pandocbounded[1]{% scales image to fit in text height/width
  \sbox\pandoc@box{#1}%
  \Gscale@div\@tempa{\textheight}{\dimexpr\ht\pandoc@box+\dp\pandoc@box\relax}%
  \Gscale@div\@tempb{\linewidth}{\wd\pandoc@box}%
  \ifdim\@tempb\p@<\@tempa\p@\let\@tempa\@tempb\fi% select the smaller of both
  \ifdim\@tempa\p@<\p@\scalebox{\@tempa}{\usebox\pandoc@box}%
  \else\usebox{\pandoc@box}%
  \fi%
}
% Set default figure placement to htbp
\def\fps@figure{htbp}
\makeatother
\setlength{\emergencystretch}{3em} % prevent overfull lines
\providecommand{\tightlist}{%
  \setlength{\itemsep}{0pt}\setlength{\parskip}{0pt}}
% Slide layout defaults for warning-clean scientific decks.
\usepackage{etoolbox}
\IfFileExists{xurl.sty}{\usepackage{xurl}}{}
\IfFileExists{seqsplit.sty}{\usepackage{seqsplit}}{\newcommand{\seqsplit}[1]{#1}}
\protected\def\breakseq#1{\seqsplit{#1}}
\protected\def\breaktt#1{\begingroup\ttfamily\seqsplit{#1}\endgroup}
\setlength{\emergencystretch}{6em}
\tolerance=5000
\hbadness=10000
\hfuzz=1pt
\setlength{\tabcolsep}{2pt}
\AtBeginEnvironment{longtable}{\tiny\renewcommand{\arraystretch}{0.86}\setlength{\tabcolsep}{1pt}}
\AtBeginEnvironment{tabular}{\tiny\renewcommand{\arraystretch}{0.86}\setlength{\tabcolsep}{1pt}}

% Natbib and cross-reference fallbacks — slides don't load natbib
% or cleveref, but combined-PDF manuscript prose may emit these
% commands. The fallback renders citations as a bracketed key list
% and cross-references as detokenized labels so slides don't fail on
% undefined control sequences or raw underscores. \providecommand is
% a no-op if packages load later.
\providecommand{\citep}[1]{[#1]}
\providecommand{\citet}[1]{#1}
\providecommand{\citealp}[1]{#1}
\providecommand{\citeauthor}[1]{#1}
\providecommand{\citeyear}[1]{#1}
\providecommand{\cref}[1]{\texttt{\detokenize{#1}}}
\providecommand{\Cref}[1]{\texttt{\detokenize{#1}}}

\newtheorem{proposition}{Proposition}
\newtheorem{hypothesis}{Hypothesis}
\usepackage{bookmark}
\IfFileExists{xurl.sty}{\usepackage{xurl}}{} % add URL line breaks if available
\urlstyle{same}
% fallback for those not using the hyperref driver hyperxmp:
\makeatletter
\@ifundefined{xmpquote}{\newcommand{\xmpquote}[1]{#1}}{}
\makeatother
\hypersetup{
  hidelinks,
  pdfcreator={LaTeX via pandoc}}

\author{\texorpdfstring{}{}}
\date{}

\begin{document}

\section{Quality validation}\label{sec:quality}

\begin{frame}[allowframebreaks]{Quality validation}
The quality gate is what separates a descriptor that merely
\emph{claims} to describe a dataset from one that \emph{provably} does.
It has two layers: structural validation of the descriptor, and
byte-level reconciliation of the descriptor against the files it names.
\end{frame}

\begin{frame}[fragile,allowframebreaks]{The validation gate reacts to
defects}
\protect\phantomsection\label{the-validation-gate-reacts-to-defects}
\breaktt{validate\_descriptor()} emits severity-tagged findings for
missing required keys, unsafe or duplicate file paths, malformed
checksums, unknown media types, duplicate or malformed fields,
constraint violations, missing primary keys, and thin provenance.
\breaktt{build\_descriptor\_report()} folds these into a readiness score,
penalising each error and warning.

To show the gate is not vacuous, \breakseq{[@fig:quality\_gate]} compares the
shipped fixture against a deliberately-perturbed copy produced by
\breaktt{demo\_broken\_descriptor()} (which strips the license, corrupts
a checksum, zeroes a row count, and removes a unit). The clean fixture
produces zero findings; the perturbed demo produces several errors and
warnings.

\begin{figure}
\centering
\pandocbounded{\includegraphics[keepaspectratio,alt={Quality gate: validation findings by severity for the clean fixture descriptor (left, zero findings) versus a deliberately-broken demonstration copy (right). The demonstration perturbation is clearly named in code and never treated as real data --- it exists only to prove the gate reacts. Counts come from severity\_counts().}]{../figures/quality_gate.png}}
\caption{Quality gate: validation findings by severity for the clean
fixture descriptor (left, zero findings) versus a deliberately-broken
demonstration copy (right). The demonstration perturbation is clearly
named in code and never treated as real data --- it exists only to prove
the gate reacts. Counts come from
\protect\breaktt{severity\_counts()}.}\label{fig:quality_gate}
\end{figure}
\end{frame}

\begin{frame}[fragile,allowframebreaks]{Descriptor↔file verification}
\protect\phantomsection\label{descriptorfile-verification}
A checksum is only meaningful if something checks it.
\breaktt{verify\_descriptor\_files()} recomputes the sha256 digest and
data-row count of every declared file that is present on disk and
compares them to the descriptor's declarations, reporting each file as
\texttt{verified}, \breaktt{checksum\_mismatch}, \texttt{row\_mismatch},
or \texttt{absent}. Files a descriptor references but does not bundle
are reported as \texttt{absent} rather than as failures --- verification
is asserted only for bytes that are actually present, which is what
makes the metadata-only release boundary honest.

\breakseq{[@fig:checksum\_verification]} shows this reconciliation for the
shipped fixture: both files verify, declared and actual row counts
agree, and each recomputed checksum matches its declaration, so the
readiness score is unpenalised.

\begin{figure}
\centering
\pandocbounded{\includegraphics[keepaspectratio,alt={Descriptor↔file verification: declared versus recomputed row counts and checksum agreement per file, from verify\_descriptor\_files(). Both shipped fixture files verify against their bytes on disk; the title reports the descriptor's readiness score. If a fixture file were edited without updating the descriptor, the corresponding row would flip to a mismatch status and the test suite would fail.}]{../figures/checksum_verification.png}}
\caption{Descriptor↔file verification: declared versus recomputed row
counts and checksum agreement per file, from
\protect\breaktt{verify\_descriptor\_files()}. Both shipped fixture files
verify against their bytes on disk; the title reports the descriptor's
readiness score. If a fixture file were edited without updating the
descriptor, the corresponding row would flip to a mismatch status and
the test suite would fail.}\label{fig:checksum_verification}
\end{figure}
\end{frame}

\begin{frame}[fragile,allowframebreaks]{Machine-readable release
manifest}
\protect\phantomsection\label{machine-readable-release-manifest}
\breaktt{build\_release\_bundle\_manifest()} packages the above into a
deterministic, JSON-ready manifest containing the schema fingerprint,
per-file checksums and row counts, field-level unit/bounds/enumeration
summaries, provenance steps, and the readiness verdict --- carrying
metadata and checksums, never the dataset bytes. The thin script
\href{../scripts/generate_release_artifacts.py}{\breaktt{scripts/generate\_release\_artifacts.py}}
writes this manifest (plus the readiness report and constraint summary)
under \texttt{output/reports/} for pre-publication review.
\end{frame}

\begin{frame}[fragile,allowframebreaks]{Validation evidence}
\protect\phantomsection\label{validation-evidence}
Every capability above is exercised by the zero-mock \texttt{tests/}
suite: structural validation and negative controls, order-independent
schema fingerprints, byte-level verification against real temporary
files, figure-data preparation, and an integration test that runs the
figure script end to end and asserts real PNGs are written. Coverage
exceeds the 90\% project gate with no mocks.
\end{frame}

\end{document}
